% ============================================================================
% _estilo.tex — sistema de diseño compartido para los PDF del frente
% Atención Armónica (Phideus). Derivado del sistema usado en el informe
% narrado de la memoria colectiva: paleta teal/dorado, cajas tcolorbox,
% tablas booktabs, más un vocabulario TikZ para los diagramas.
%
% Uso:  % ============================================================================
% _estilo.tex — sistema de diseño compartido para los PDF del frente
% Atención Armónica (Phideus). Derivado del sistema usado en el informe
% narrado de la memoria colectiva: paleta teal/dorado, cajas tcolorbox,
% tablas booktabs, más un vocabulario TikZ para los diagramas.
%
% Uso:  % ============================================================================
% _estilo.tex — sistema de diseño compartido para los PDF del frente
% Atención Armónica (Phideus). Derivado del sistema usado en el informe
% narrado de la memoria colectiva: paleta teal/dorado, cajas tcolorbox,
% tablas booktabs, más un vocabulario TikZ para los diagramas.
%
% Uso:  % ============================================================================
% _estilo.tex — sistema de diseño compartido para los PDF del frente
% Atención Armónica (Phideus). Derivado del sistema usado en el informe
% narrado de la memoria colectiva: paleta teal/dorado, cajas tcolorbox,
% tablas booktabs, más un vocabulario TikZ para los diagramas.
%
% Uso:  \input{_estilo}  justo antes de \begin{document}
%       \titleblock{Título}{Subtítulo}
% Compilar dos veces:  pdflatex -interaction=nonstopmode NOMBRE.tex
% ============================================================================

\usepackage[utf8]{inputenc}
\usepackage[T1]{fontenc}
\usepackage[spanish,es-noshorthands]{babel}
\usepackage{lmodern}
\usepackage[a4paper,margin=2.3cm]{geometry}
\usepackage{microtype}

\usepackage{xcolor}
\usepackage{amsmath}
\usepackage{amssymb}
\usepackage{booktabs}
\usepackage{colortbl}
\usepackage{array}
\usepackage{tabularx}
\usepackage{float}
\usepackage[most]{tcolorbox}
\usepackage{titlesec}
\usepackage{fancyhdr}
\usepackage{enumitem}
\usepackage{tikz}
\usetikzlibrary{positioning,arrows.meta,fit,backgrounds,calc,shapes.geometric,shapes.misc}
\usepackage[hidelinks]{hyperref}

% ---- Paleta ---------------------------------------------------------------
\definecolor{accent}{RGB}{20,105,95}      % teal profundo
\definecolor{accent2}{RGB}{196,142,53}    % dorado
\definecolor{inkgray}{RGB}{90,90,90}
\definecolor{rowgray}{RGB}{245,247,246}
\definecolor{alertred}{RGB}{160,58,48}

% ---- Secciones ------------------------------------------------------------
\titleformat{\section}
  {\sffamily\Large\bfseries\color{accent}}{\thesection}{0.6em}{}
  [{\color{accent2}\titlerule[1.4pt]}]
\titleformat{\subsection}
  {\sffamily\large\bfseries\color{accent}}{\thesubsection}{0.5em}{}
\titlespacing*{\section}{0pt}{1.4\baselineskip}{0.6\baselineskip}

% ---- Pie ------------------------------------------------------------------
\newcommand{\footername}{}
\pagestyle{fancy}
\fancyhf{}
\renewcommand{\headrulewidth}{0pt}
\renewcommand{\footrulewidth}{0.4pt}
\renewcommand{\footrule}{\hbox to\headwidth{\color{accent}\leaders\hrule height \footrulewidth\hfill}}
\fancyfoot[L]{\footnotesize\color{inkgray}\sffamily \footername}
\fancyfoot[R]{\footnotesize\color{inkgray}\sffamily \thepage}

% ---- Cajas ----------------------------------------------------------------
\newtcolorbox{keybox}[1]{
  enhanced, breakable, colback=accent!7, colframe=accent,
  boxrule=0.8pt, arc=2pt, left=8pt, right=8pt, top=6pt, bottom=6pt,
  coltitle=white, fonttitle=\sffamily\bfseries, title={#1}}

\newtcolorbox{alertbox}[1]{
  enhanced, breakable, colback=alertred!6, colframe=alertred,
  boxrule=0.8pt, arc=2pt, left=8pt, right=8pt, top=6pt, bottom=6pt,
  coltitle=white, fonttitle=\sffamily\bfseries, title={#1}}

\newtcolorbox{quietbox}{
  enhanced, breakable, colback=rowgray, colframe=rowgray,
  boxrule=0pt, arc=1pt, left=8pt, right=8pt, top=5pt, bottom=5pt,
  borderline west={2pt}{0pt}{accent!55}}

\newtcolorbox{goldbox}{
  enhanced, breakable, colback=accent2!8, colframe=accent2!8,
  boxrule=0pt, arc=1pt, left=8pt, right=8pt, top=5pt, bottom=5pt,
  borderline west={2pt}{0pt}{accent2}}

\newcommand{\headrow}{\rowcolor{accent!12}}
\setlist{leftmargin=1.4em, topsep=2pt, itemsep=1pt}
\renewcommand{\arraystretch}{1.15}

% ---- Bloque de título ------------------------------------------------------
\newcommand{\titleblock}[2]{%
  \noindent{\color{accent2}\rule{\linewidth}{1.6pt}}\par\smallskip
  {\sffamily\bfseries\Huge\color{accent} #1\par}\smallskip
  \ifx\relax#2\relax\else{\sffamily\large\color{inkgray} #2\par}\fi
  \smallskip\noindent{\color{accent2}\rule{\linewidth}{1.6pt}}\par\bigskip
  \gdef\footername{#1}%
}

% ---- Vocabulario de figuras TikZ ------------------------------------------
% Los diagramas ASCII de los originales se redibujan con estos estilos.
\tikzset{
  fbox/.style   ={draw=accent!65, fill=accent!8,  rounded corners=3pt, align=center,
                  inner sep=6pt, font=\sffamily\small, minimum height=8mm},
  fgold/.style  ={draw=accent2!85, fill=accent2!12, rounded corners=3pt, align=center,
                  inner sep=6pt, font=\sffamily\small, minimum height=8mm},
  fbad/.style   ={draw=alertred!70, fill=alertred!7, rounded corners=3pt, align=center,
                  inner sep=6pt, font=\sffamily\small, minimum height=8mm},
  fplain/.style ={draw=inkgray!45, fill=rowgray, rounded corners=3pt, align=center,
                  inner sep=6pt, font=\sffamily\small, minimum height=8mm},
  fdot/.style   ={circle, draw=accent!70, fill=accent!12, inner sep=0pt,
                  minimum size=6.5mm, font=\sffamily\scriptsize},
  fdot2/.style  ={circle, draw=accent2!90, fill=accent2!15, inner sep=0pt,
                  minimum size=6.5mm, font=\sffamily\scriptsize},
  flab/.style   ={font=\sffamily\scriptsize\color{inkgray}, align=center, inner sep=2pt},
  farr/.style   ={-{Stealth[length=2.6mm]}, draw=accent, thick},
  fgarr/.style  ={-{Stealth[length=2.6mm]}, draw=accent2!90!black, thick},
  fbarr/.style  ={-{Stealth[length=2.6mm]}, draw=alertred, thick},
  fdash/.style  ={-{Stealth[length=2.6mm]}, draw=inkgray!70, thick, dashed},
  fline/.style  ={draw=accent!55, thick},
  fxline/.style ={draw=alertred!80, thick, dash pattern=on 3pt off 2pt},
}

% Pie de figura, sin depender del paquete caption
\newcommand{\figcap}[1]{\par\smallskip{\sffamily\footnotesize\color{inkgray} #1\par}}

% Caja monoespaciada — SOLO para contenido ASCII puro (sin dibujo de cajas)
\newtcolorbox{codebox}{enhanced, breakable, colback=rowgray, colframe=accent!25,
  boxrule=0.5pt, arc=2pt, left=6pt, right=6pt, top=4pt, bottom=4pt,
  fontupper=\ttfamily\scriptsize}
  justo antes de \begin{document}
%       \titleblock{Título}{Subtítulo}
% Compilar dos veces:  pdflatex -interaction=nonstopmode NOMBRE.tex
% ============================================================================

\usepackage[utf8]{inputenc}
\usepackage[T1]{fontenc}
\usepackage[spanish,es-noshorthands]{babel}
\usepackage{lmodern}
\usepackage[a4paper,margin=2.3cm]{geometry}
\usepackage{microtype}

\usepackage{xcolor}
\usepackage{amsmath}
\usepackage{amssymb}
\usepackage{booktabs}
\usepackage{colortbl}
\usepackage{array}
\usepackage{tabularx}
\usepackage{float}
\usepackage[most]{tcolorbox}
\usepackage{titlesec}
\usepackage{fancyhdr}
\usepackage{enumitem}
\usepackage{tikz}
\usetikzlibrary{positioning,arrows.meta,fit,backgrounds,calc,shapes.geometric,shapes.misc}
\usepackage[hidelinks]{hyperref}

% ---- Paleta ---------------------------------------------------------------
\definecolor{accent}{RGB}{20,105,95}      % teal profundo
\definecolor{accent2}{RGB}{196,142,53}    % dorado
\definecolor{inkgray}{RGB}{90,90,90}
\definecolor{rowgray}{RGB}{245,247,246}
\definecolor{alertred}{RGB}{160,58,48}

% ---- Secciones ------------------------------------------------------------
\titleformat{\section}
  {\sffamily\Large\bfseries\color{accent}}{\thesection}{0.6em}{}
  [{\color{accent2}\titlerule[1.4pt]}]
\titleformat{\subsection}
  {\sffamily\large\bfseries\color{accent}}{\thesubsection}{0.5em}{}
\titlespacing*{\section}{0pt}{1.4\baselineskip}{0.6\baselineskip}

% ---- Pie ------------------------------------------------------------------
\newcommand{\footername}{}
\pagestyle{fancy}
\fancyhf{}
\renewcommand{\headrulewidth}{0pt}
\renewcommand{\footrulewidth}{0.4pt}
\renewcommand{\footrule}{\hbox to\headwidth{\color{accent}\leaders\hrule height \footrulewidth\hfill}}
\fancyfoot[L]{\footnotesize\color{inkgray}\sffamily \footername}
\fancyfoot[R]{\footnotesize\color{inkgray}\sffamily \thepage}

% ---- Cajas ----------------------------------------------------------------
\newtcolorbox{keybox}[1]{
  enhanced, breakable, colback=accent!7, colframe=accent,
  boxrule=0.8pt, arc=2pt, left=8pt, right=8pt, top=6pt, bottom=6pt,
  coltitle=white, fonttitle=\sffamily\bfseries, title={#1}}

\newtcolorbox{alertbox}[1]{
  enhanced, breakable, colback=alertred!6, colframe=alertred,
  boxrule=0.8pt, arc=2pt, left=8pt, right=8pt, top=6pt, bottom=6pt,
  coltitle=white, fonttitle=\sffamily\bfseries, title={#1}}

\newtcolorbox{quietbox}{
  enhanced, breakable, colback=rowgray, colframe=rowgray,
  boxrule=0pt, arc=1pt, left=8pt, right=8pt, top=5pt, bottom=5pt,
  borderline west={2pt}{0pt}{accent!55}}

\newtcolorbox{goldbox}{
  enhanced, breakable, colback=accent2!8, colframe=accent2!8,
  boxrule=0pt, arc=1pt, left=8pt, right=8pt, top=5pt, bottom=5pt,
  borderline west={2pt}{0pt}{accent2}}

\newcommand{\headrow}{\rowcolor{accent!12}}
\setlist{leftmargin=1.4em, topsep=2pt, itemsep=1pt}
\renewcommand{\arraystretch}{1.15}

% ---- Bloque de título ------------------------------------------------------
\newcommand{\titleblock}[2]{%
  \noindent{\color{accent2}\rule{\linewidth}{1.6pt}}\par\smallskip
  {\sffamily\bfseries\Huge\color{accent} #1\par}\smallskip
  \ifx\relax#2\relax\else{\sffamily\large\color{inkgray} #2\par}\fi
  \smallskip\noindent{\color{accent2}\rule{\linewidth}{1.6pt}}\par\bigskip
  \gdef\footername{#1}%
}

% ---- Vocabulario de figuras TikZ ------------------------------------------
% Los diagramas ASCII de los originales se redibujan con estos estilos.
\tikzset{
  fbox/.style   ={draw=accent!65, fill=accent!8,  rounded corners=3pt, align=center,
                  inner sep=6pt, font=\sffamily\small, minimum height=8mm},
  fgold/.style  ={draw=accent2!85, fill=accent2!12, rounded corners=3pt, align=center,
                  inner sep=6pt, font=\sffamily\small, minimum height=8mm},
  fbad/.style   ={draw=alertred!70, fill=alertred!7, rounded corners=3pt, align=center,
                  inner sep=6pt, font=\sffamily\small, minimum height=8mm},
  fplain/.style ={draw=inkgray!45, fill=rowgray, rounded corners=3pt, align=center,
                  inner sep=6pt, font=\sffamily\small, minimum height=8mm},
  fdot/.style   ={circle, draw=accent!70, fill=accent!12, inner sep=0pt,
                  minimum size=6.5mm, font=\sffamily\scriptsize},
  fdot2/.style  ={circle, draw=accent2!90, fill=accent2!15, inner sep=0pt,
                  minimum size=6.5mm, font=\sffamily\scriptsize},
  flab/.style   ={font=\sffamily\scriptsize\color{inkgray}, align=center, inner sep=2pt},
  farr/.style   ={-{Stealth[length=2.6mm]}, draw=accent, thick},
  fgarr/.style  ={-{Stealth[length=2.6mm]}, draw=accent2!90!black, thick},
  fbarr/.style  ={-{Stealth[length=2.6mm]}, draw=alertred, thick},
  fdash/.style  ={-{Stealth[length=2.6mm]}, draw=inkgray!70, thick, dashed},
  fline/.style  ={draw=accent!55, thick},
  fxline/.style ={draw=alertred!80, thick, dash pattern=on 3pt off 2pt},
}

% Pie de figura, sin depender del paquete caption
\newcommand{\figcap}[1]{\par\smallskip{\sffamily\footnotesize\color{inkgray} #1\par}}

% Caja monoespaciada — SOLO para contenido ASCII puro (sin dibujo de cajas)
\newtcolorbox{codebox}{enhanced, breakable, colback=rowgray, colframe=accent!25,
  boxrule=0.5pt, arc=2pt, left=6pt, right=6pt, top=4pt, bottom=4pt,
  fontupper=\ttfamily\scriptsize}
  justo antes de \begin{document}
%       \titleblock{Título}{Subtítulo}
% Compilar dos veces:  pdflatex -interaction=nonstopmode NOMBRE.tex
% ============================================================================

\usepackage[utf8]{inputenc}
\usepackage[T1]{fontenc}
\usepackage[spanish,es-noshorthands]{babel}
\usepackage{lmodern}
\usepackage[a4paper,margin=2.3cm]{geometry}
\usepackage{microtype}

\usepackage{xcolor}
\usepackage{amsmath}
\usepackage{amssymb}
\usepackage{booktabs}
\usepackage{colortbl}
\usepackage{array}
\usepackage{tabularx}
\usepackage{float}
\usepackage[most]{tcolorbox}
\usepackage{titlesec}
\usepackage{fancyhdr}
\usepackage{enumitem}
\usepackage{tikz}
\usetikzlibrary{positioning,arrows.meta,fit,backgrounds,calc,shapes.geometric,shapes.misc}
\usepackage[hidelinks]{hyperref}

% ---- Paleta ---------------------------------------------------------------
\definecolor{accent}{RGB}{20,105,95}      % teal profundo
\definecolor{accent2}{RGB}{196,142,53}    % dorado
\definecolor{inkgray}{RGB}{90,90,90}
\definecolor{rowgray}{RGB}{245,247,246}
\definecolor{alertred}{RGB}{160,58,48}

% ---- Secciones ------------------------------------------------------------
\titleformat{\section}
  {\sffamily\Large\bfseries\color{accent}}{\thesection}{0.6em}{}
  [{\color{accent2}\titlerule[1.4pt]}]
\titleformat{\subsection}
  {\sffamily\large\bfseries\color{accent}}{\thesubsection}{0.5em}{}
\titlespacing*{\section}{0pt}{1.4\baselineskip}{0.6\baselineskip}

% ---- Pie ------------------------------------------------------------------
\newcommand{\footername}{}
\pagestyle{fancy}
\fancyhf{}
\renewcommand{\headrulewidth}{0pt}
\renewcommand{\footrulewidth}{0.4pt}
\renewcommand{\footrule}{\hbox to\headwidth{\color{accent}\leaders\hrule height \footrulewidth\hfill}}
\fancyfoot[L]{\footnotesize\color{inkgray}\sffamily \footername}
\fancyfoot[R]{\footnotesize\color{inkgray}\sffamily \thepage}

% ---- Cajas ----------------------------------------------------------------
\newtcolorbox{keybox}[1]{
  enhanced, breakable, colback=accent!7, colframe=accent,
  boxrule=0.8pt, arc=2pt, left=8pt, right=8pt, top=6pt, bottom=6pt,
  coltitle=white, fonttitle=\sffamily\bfseries, title={#1}}

\newtcolorbox{alertbox}[1]{
  enhanced, breakable, colback=alertred!6, colframe=alertred,
  boxrule=0.8pt, arc=2pt, left=8pt, right=8pt, top=6pt, bottom=6pt,
  coltitle=white, fonttitle=\sffamily\bfseries, title={#1}}

\newtcolorbox{quietbox}{
  enhanced, breakable, colback=rowgray, colframe=rowgray,
  boxrule=0pt, arc=1pt, left=8pt, right=8pt, top=5pt, bottom=5pt,
  borderline west={2pt}{0pt}{accent!55}}

\newtcolorbox{goldbox}{
  enhanced, breakable, colback=accent2!8, colframe=accent2!8,
  boxrule=0pt, arc=1pt, left=8pt, right=8pt, top=5pt, bottom=5pt,
  borderline west={2pt}{0pt}{accent2}}

\newcommand{\headrow}{\rowcolor{accent!12}}
\setlist{leftmargin=1.4em, topsep=2pt, itemsep=1pt}
\renewcommand{\arraystretch}{1.15}

% ---- Bloque de título ------------------------------------------------------
\newcommand{\titleblock}[2]{%
  \noindent{\color{accent2}\rule{\linewidth}{1.6pt}}\par\smallskip
  {\sffamily\bfseries\Huge\color{accent} #1\par}\smallskip
  \ifx\relax#2\relax\else{\sffamily\large\color{inkgray} #2\par}\fi
  \smallskip\noindent{\color{accent2}\rule{\linewidth}{1.6pt}}\par\bigskip
  \gdef\footername{#1}%
}

% ---- Vocabulario de figuras TikZ ------------------------------------------
% Los diagramas ASCII de los originales se redibujan con estos estilos.
\tikzset{
  fbox/.style   ={draw=accent!65, fill=accent!8,  rounded corners=3pt, align=center,
                  inner sep=6pt, font=\sffamily\small, minimum height=8mm},
  fgold/.style  ={draw=accent2!85, fill=accent2!12, rounded corners=3pt, align=center,
                  inner sep=6pt, font=\sffamily\small, minimum height=8mm},
  fbad/.style   ={draw=alertred!70, fill=alertred!7, rounded corners=3pt, align=center,
                  inner sep=6pt, font=\sffamily\small, minimum height=8mm},
  fplain/.style ={draw=inkgray!45, fill=rowgray, rounded corners=3pt, align=center,
                  inner sep=6pt, font=\sffamily\small, minimum height=8mm},
  fdot/.style   ={circle, draw=accent!70, fill=accent!12, inner sep=0pt,
                  minimum size=6.5mm, font=\sffamily\scriptsize},
  fdot2/.style  ={circle, draw=accent2!90, fill=accent2!15, inner sep=0pt,
                  minimum size=6.5mm, font=\sffamily\scriptsize},
  flab/.style   ={font=\sffamily\scriptsize\color{inkgray}, align=center, inner sep=2pt},
  farr/.style   ={-{Stealth[length=2.6mm]}, draw=accent, thick},
  fgarr/.style  ={-{Stealth[length=2.6mm]}, draw=accent2!90!black, thick},
  fbarr/.style  ={-{Stealth[length=2.6mm]}, draw=alertred, thick},
  fdash/.style  ={-{Stealth[length=2.6mm]}, draw=inkgray!70, thick, dashed},
  fline/.style  ={draw=accent!55, thick},
  fxline/.style ={draw=alertred!80, thick, dash pattern=on 3pt off 2pt},
}

% Pie de figura, sin depender del paquete caption
\newcommand{\figcap}[1]{\par\smallskip{\sffamily\footnotesize\color{inkgray} #1\par}}

% Caja monoespaciada — SOLO para contenido ASCII puro (sin dibujo de cajas)
\newtcolorbox{codebox}{enhanced, breakable, colback=rowgray, colframe=accent!25,
  boxrule=0.5pt, arc=2pt, left=6pt, right=6pt, top=4pt, bottom=4pt,
  fontupper=\ttfamily\scriptsize}
  justo antes de \begin{document}
%       \titleblock{Título}{Subtítulo}
% Compilar dos veces:  pdflatex -interaction=nonstopmode NOMBRE.tex
% ============================================================================

\usepackage[utf8]{inputenc}
\usepackage[T1]{fontenc}
\usepackage[spanish,es-noshorthands]{babel}
\usepackage{lmodern}
\usepackage[a4paper,margin=2.3cm]{geometry}
\usepackage{microtype}

\usepackage{xcolor}
\usepackage{amsmath}
\usepackage{amssymb}
\usepackage{booktabs}
\usepackage{colortbl}
\usepackage{array}
\usepackage{tabularx}
\usepackage{float}
\usepackage[most]{tcolorbox}
\usepackage{titlesec}
\usepackage{fancyhdr}
\usepackage{enumitem}
\usepackage{tikz}
\usetikzlibrary{positioning,arrows.meta,fit,backgrounds,calc,shapes.geometric,shapes.misc}
\usepackage[hidelinks]{hyperref}

% ---- Paleta ---------------------------------------------------------------
\definecolor{accent}{RGB}{20,105,95}      % teal profundo
\definecolor{accent2}{RGB}{196,142,53}    % dorado
\definecolor{inkgray}{RGB}{90,90,90}
\definecolor{rowgray}{RGB}{245,247,246}
\definecolor{alertred}{RGB}{160,58,48}

% ---- Secciones ------------------------------------------------------------
\titleformat{\section}
  {\sffamily\Large\bfseries\color{accent}}{\thesection}{0.6em}{}
  [{\color{accent2}\titlerule[1.4pt]}]
\titleformat{\subsection}
  {\sffamily\large\bfseries\color{accent}}{\thesubsection}{0.5em}{}
\titlespacing*{\section}{0pt}{1.4\baselineskip}{0.6\baselineskip}

% ---- Pie ------------------------------------------------------------------
\newcommand{\footername}{}
\pagestyle{fancy}
\fancyhf{}
\renewcommand{\headrulewidth}{0pt}
\renewcommand{\footrulewidth}{0.4pt}
\renewcommand{\footrule}{\hbox to\headwidth{\color{accent}\leaders\hrule height \footrulewidth\hfill}}
\fancyfoot[L]{\footnotesize\color{inkgray}\sffamily \footername}
\fancyfoot[R]{\footnotesize\color{inkgray}\sffamily \thepage}

% ---- Cajas ----------------------------------------------------------------
\newtcolorbox{keybox}[1]{
  enhanced, breakable, colback=accent!7, colframe=accent,
  boxrule=0.8pt, arc=2pt, left=8pt, right=8pt, top=6pt, bottom=6pt,
  coltitle=white, fonttitle=\sffamily\bfseries, title={#1}}

\newtcolorbox{alertbox}[1]{
  enhanced, breakable, colback=alertred!6, colframe=alertred,
  boxrule=0.8pt, arc=2pt, left=8pt, right=8pt, top=6pt, bottom=6pt,
  coltitle=white, fonttitle=\sffamily\bfseries, title={#1}}

\newtcolorbox{quietbox}{
  enhanced, breakable, colback=rowgray, colframe=rowgray,
  boxrule=0pt, arc=1pt, left=8pt, right=8pt, top=5pt, bottom=5pt,
  borderline west={2pt}{0pt}{accent!55}}

\newtcolorbox{goldbox}{
  enhanced, breakable, colback=accent2!8, colframe=accent2!8,
  boxrule=0pt, arc=1pt, left=8pt, right=8pt, top=5pt, bottom=5pt,
  borderline west={2pt}{0pt}{accent2}}

\newcommand{\headrow}{\rowcolor{accent!12}}
\setlist{leftmargin=1.4em, topsep=2pt, itemsep=1pt}
\renewcommand{\arraystretch}{1.15}

% ---- Bloque de título ------------------------------------------------------
\newcommand{\titleblock}[2]{%
  \noindent{\color{accent2}\rule{\linewidth}{1.6pt}}\par\smallskip
  {\sffamily\bfseries\Huge\color{accent} #1\par}\smallskip
  \ifx\relax#2\relax\else{\sffamily\large\color{inkgray} #2\par}\fi
  \smallskip\noindent{\color{accent2}\rule{\linewidth}{1.6pt}}\par\bigskip
  \gdef\footername{#1}%
}

% ---- Vocabulario de figuras TikZ ------------------------------------------
% Los diagramas ASCII de los originales se redibujan con estos estilos.
\tikzset{
  fbox/.style   ={draw=accent!65, fill=accent!8,  rounded corners=3pt, align=center,
                  inner sep=6pt, font=\sffamily\small, minimum height=8mm},
  fgold/.style  ={draw=accent2!85, fill=accent2!12, rounded corners=3pt, align=center,
                  inner sep=6pt, font=\sffamily\small, minimum height=8mm},
  fbad/.style   ={draw=alertred!70, fill=alertred!7, rounded corners=3pt, align=center,
                  inner sep=6pt, font=\sffamily\small, minimum height=8mm},
  fplain/.style ={draw=inkgray!45, fill=rowgray, rounded corners=3pt, align=center,
                  inner sep=6pt, font=\sffamily\small, minimum height=8mm},
  fdot/.style   ={circle, draw=accent!70, fill=accent!12, inner sep=0pt,
                  minimum size=6.5mm, font=\sffamily\scriptsize},
  fdot2/.style  ={circle, draw=accent2!90, fill=accent2!15, inner sep=0pt,
                  minimum size=6.5mm, font=\sffamily\scriptsize},
  flab/.style   ={font=\sffamily\scriptsize\color{inkgray}, align=center, inner sep=2pt},
  farr/.style   ={-{Stealth[length=2.6mm]}, draw=accent, thick},
  fgarr/.style  ={-{Stealth[length=2.6mm]}, draw=accent2!90!black, thick},
  fbarr/.style  ={-{Stealth[length=2.6mm]}, draw=alertred, thick},
  fdash/.style  ={-{Stealth[length=2.6mm]}, draw=inkgray!70, thick, dashed},
  fline/.style  ={draw=accent!55, thick},
  fxline/.style ={draw=alertred!80, thick, dash pattern=on 3pt off 2pt},
}

% Pie de figura, sin depender del paquete caption
\newcommand{\figcap}[1]{\par\smallskip{\sffamily\footnotesize\color{inkgray} #1\par}}

% Caja monoespaciada — SOLO para contenido ASCII puro (sin dibujo de cajas)
\newtcolorbox{codebox}{enhanced, breakable, colback=rowgray, colframe=accent!25,
  boxrule=0.5pt, arc=2pt, left=6pt, right=6pt, top=4pt, bottom=4pt,
  fontupper=\ttfamily\scriptsize}
